\chapter{Conception de la solution}

\section*{Introduction}
\addcontentsline{toc}{section}{Introduction}

Ce chapitre décrit la conception architecturale du BTE Security AI Agent. Il part de la
vue globale de l'architecture logicielle, de ses onze conteneurs et de son point
d'entrée, pour entrer ensuite dans la conception de chacun des six sous-systèmes qui
composent la solution : l'inférence LLM, la revue de \textit{Pull Request}, l'analyse du
diff, l'assistant de chat opérationnel, la persistance et l'observabilité.

% ─────────────────────────────────────────────────────────────────────────────
\section{Présentation de l'architecture logicielle}

    \subsection{Architecture en couches}

        La plateforme BTE Security AI Agent est déployée sous forme de onze conteneurs
        Docker organisés en quatre couches fonctionnelles complémentaires. La figure
        \ref{fig:full_arch} présente cette architecture globale avec l'ensemble des
        interactions entre les composants.

        \begin{figure}[H]
            \centering
            \frame{\includegraphics[width=\textwidth]{full_architecture}}
            \caption{Architecture globale du BTE Security AI Agent}
            \label{fig:full_arch}
        \end{figure}

        La \textbf{couche entrée} est assurée par nginx \cite{NGINX} en tant que reverse
        proxy et unique point d'entrée externe sur les ports 80 et 443. Il achemine les
        requêtes vers les services internes selon le chemin URL et applique une
        authentification HTTP Basic sur les interfaces sensibles. La \textbf{couche IA}
        comprend le FastAPI \cite{FastAPI} agent, qui orchestre le workflow LangGraph
        \cite{LangGraph}, appelle les modèles LLM et exécute les scanners, Ollama
        \cite{Ollama}, le serveur d'inférence LLM accessible uniquement depuis le réseau
        Docker \cite{Docker} interne, et Open WebUI qui offre une interface graphique
        complémentaire pour interagir directement avec les modèles Ollama lors des phases
        de mise au point. La \textbf{couche données} est constituée de
        PostgreSQL \cite{PostgreSQL}, qui héberge la base de connaissances sécurité et les
        checkpoints LangGraph, et de Redis \cite{Redis}, qui assure la déduplication des
        webhooks, le rate limiting et le cache des résultats de scanners. Enfin, la
        \textbf{couche observabilité} regroupe Prometheus \cite{Prometheus},
        VictoriaMetrics \cite{VictoriaMetrics}, Grafana \cite{Grafana},
        node-exporter \cite{NodeExporter} et AlertManager \cite{AlertManager}, qui
        assurent la collecte, le stockage, la visualisation et l'alerting.

    \subsection{Inventaire des conteneurs}

        Le tableau \ref{tab:containers} présente l'inventaire détaillé des onze conteneurs
        avec leurs caractéristiques principales.

        \begin{table}[H]
        \centering
        \begin{tabular}{|m{3.5cm}|m{3cm}|m{2cm}|m{2cm}|m{2.5cm}|}
        \hline
        \textbf{Conteneur} & \textbf{Image} & \textbf{Port} & \textbf{RAM} & \textbf{Rôle} \\
        \hline
        \texttt{devsecops-agent} & Python 3.12 custom & 8000 & \textasciitilde120 Mo & Cœur IA + API \\
        \hline
        \texttt{ollama} & \texttt{ollama/\allowbreak{}ollama:latest} & interne & \textasciitilde10 Go & Inférence LLM \\
        \hline
        \texttt{postgres} & \texttt{postgres:16-alpine} & interne & \textasciitilde33 Mo & BD + checkpoints \\
        \hline
        \texttt{redis} & \texttt{redis:7-alpine} & interne & \textasciitilde5 Mo & Cache + dédup \\
        \hline
        \texttt{nginx} & \texttt{nginx:alpine} & 80, 443 & \textasciitilde3 Mo & Reverse proxy \\
        \hline
        \texttt{prometheus} & \texttt{prom/\allowbreak{}prometheus} & interne & \textasciitilde26 Mo & Métriques \\
        \hline
        \texttt{alertmanager} & \texttt{prom/\allowbreak{}alertmanager} & interne & \textasciitilde10 Mo & Routage alertes \\
        \hline
        \texttt{grafana} & \texttt{grafana/\allowbreak{}grafana} & interne & \textasciitilde96 Mo & Dashboards \\
        \hline
        \texttt{victoriametrics} & \texttt{victoria-metrics} & interne & \textasciitilde97 Mo & Stockage 90j \\
        \hline
        \texttt{node-exporter} & \texttt{prom/\allowbreak{}node-exporter} & \texttt{host:9100} & \textasciitilde10 Mo & Métriques host \\
        \hline
        \texttt{open-webui} & \texttt{open-webui:main} & 3001 & \textasciitilde300 Mo & UI Ollama \\
        \hline
        \end{tabular}
        \caption{Inventaire des 11 conteneurs — ports, mémoire et rôles}
        \label{tab:containers}
        \end{table}

    \subsection{Routage et point d'entrée}

        Nginx est le point d'entrée unique de la plateforme. Le tableau
        \ref{tab:nginx_routing} présente l'ensemble des règles de routage configurées.

        \begin{table}[H]
        \centering
        \begin{tabular}{|m{3.5cm}|m{3.5cm}|m{2cm}|m{2cm}|}
        \hline
        \textbf{Location} & \textbf{Upstream} & \textbf{Timeout} & \textbf{Auth} \\
        \hline
        \texttt{GET /} & redirect → \texttt{/ui} & — & — \\
        \hline
        \texttt{GET /ui} & \texttt{agent:8000/ui} & défaut & Basic Auth \\
        \hline
        \texttt{POST /chat/} & \texttt{agent:8000/chat/} & 1800s & Basic Auth \\
        \hline
        \texttt{POST /\allowbreak{}webhooks/\allowbreak{}github} & \texttt{agent:8000} & 30s & HMAC \\
        \hline
        \texttt{/grafana/} & \texttt{grafana:3000} & défaut & Grafana auth \\
        \hline
        \texttt{/\allowbreak{}prometheus/} & \texttt{prometheus:9090} & défaut & — \\
        \hline
        \end{tabular}
        \caption{Table de routage nginx}
        \label{tab:nginx_routing}
        \end{table}

% ─────────────────────────────────────────────────────────────────────────────
\section{Description de la conception du logiciel}

    \subsection{Sous-système d'inférence LLM}

        \subsubsection{Architecture logique}

        La séparation du pipeline LLM en deux modèles distincts répond à une contrainte de
        performance. Utiliser le modèle 14B pour la classification de la
        \textit{Pull Request} serait disproportionné : cette tâche ne nécessite que la
        production d'une simple étiquette JSON en environ trente secondes, alors que la revue de
        sécurité approfondie requiert un modèle plus capable. La figure
        \ref{fig:two_model_arch} illustre cette architecture.

        \begin{figure}[H]
            \centering
            \frame{\includegraphics[width=0.9\columnwidth]{two_model_architecture}}
            \caption{Architecture à deux modèles LLM}
            \label{fig:two_model_arch}
        \end{figure}

        Les instances LLM sont construites une seule fois par processus, puis réutilisées
        pour les appels suivants. La configuration \texttt{OLLAMA\_\allowbreak{}MAX\_\allowbreak{}LOADED\_\allowbreak{}MODELS=1}
        garantit en outre qu'un seul modèle est chargé en RAM à la fois, ce qui réserve
        l'ensemble des ressources au modèle actif.

        La fusion de deux appels 14B séquentiels en un seul appel combiné dans le nœud
        \texttt{analyze\_\allowbreak{}review\_\allowbreak{}node} est l'optimisation principale du pipeline. Avant
        cette optimisation, l'analyse de sécurité et la revue de qualité de code nécessitaient
        deux appels successifs totalisant entre trente et cinquante minutes. Après la fusion,
        un seul appel produit les deux résultats en quinze à vingt-cinq minutes
        selon la complexité de la \textit{Pull Request} (mesure empirique sur CPU
        Haswell 12 cœurs sans GPU), soit un gain de près de cinquante pourcent
        sur la durée totale.

        \paragraph{Contrainte matérielle.}
        La durée d'analyse de 15 à 25 minutes par \textit{Pull Request} découle directement
        du choix d'inférence locale sur CPU sans GPU, imposé par les exigences de
        confidentialité de la BCT. Sur ce matériel, le modèle 14B produit environ 3 tokens
        par seconde. Une migration vers un GPU dédié, qui resterait sur l'infrastructure BTE
        et préserverait donc cette confidentialité, ramènerait l'analyse à quelques minutes
        par \textit{Pull Request} ; cette perspective est développée dans la conclusion
        générale.

        \subsubsection{Configuration des fabriques LLM}

        Le tableau \ref{tab:llm_factories} détaille la configuration de chaque fonction de
        fabrique LLM utilisée dans le pipeline.

        \begin{table}[H]
        \centering
        \begin{tabular}{|m{3.5cm}|m{2cm}|m{2cm}|m{2cm}|m{2cm}|m{2.5cm}|}
        \hline
        \textbf{Fonction} & \textbf{Modèle} & \textbf{num\_ctx} & \textbf{num\_predict} & \textbf{temp.} & \textbf{Durée} \\
        \hline
        \texttt{get\_\allowbreak{}fast\_\allowbreak{}llm()} & 7b & 4 096 & 512 & 0.0 & \textasciitilde30s \\
        \hline
        \texttt{get\_\allowbreak{}combined\_\allowbreak{}llm()} & 14b & 12 288 & 2 500 & 0.1 & 15–25 min \\
        \hline
        \end{tabular}
        \caption{Configuration des fonctions de fabrique LLM}
        \label{tab:llm_factories}
        \end{table}

    \subsection{Sous-système de revue de Pull Request}

        \subsubsection{Architecture logique}

        Le workflow de revue de \textit{Pull Request} est modélisé comme un
        \texttt{StateGraph} LangGraph \cite{LangGraph} avec neuf nœuds et un état partagé
        de type TypedDict nommé \texttt{PRReviewState}. La figure \ref{fig:langgraph_graph}
        présente la topologie complète du graphe.

        \begin{figure}[H]
            \centering
            \frame{\includegraphics[width=\columnwidth,height=0.80\textheight,keepaspectratio]{langgraph_state_graph}}
            \caption{Graphe d'état LangGraph}
            \label{fig:langgraph_graph}
        \end{figure}

        L'\texttt{intake\_node} est le point d'entrée du workflow. Il assure la
        déduplication Redis avec l'opération \texttt{SET NX}, le rate limiting, le clone
        shallow du dépôt avec token injecté dans l'URL, la génération du diff local avec
        quinze lignes de contexte via \texttt{git diff -U15}, et la récupération des
        dix dernières revues du dépôt depuis PostgreSQL pour injection dans le prompt. Le
        \texttt{classify\_\allowbreak{}node} appelle ensuite le modèle 7B avec \texttt{format="json"}
        et \texttt{temperature=0.\allowbreak{}0} pour classifier la PR en cinq catégories : feature,
        dependency, infrastructure, config et docs. En cas d'indisponibilité du LLM, le
        circuit breaker active le fallback regex \texttt{\_\allowbreak{}fallback\_\allowbreak{}classify()} basé sur
        les extensions de fichiers.

        Les nœuds de scanning, à savoir \texttt{scan\_\allowbreak{}full\_\allowbreak{}node} et \texttt{scan\_\allowbreak{}fs\_\allowbreak{}node},
        exécutent les scanners applicables en parallèle via \texttt{asyncio.\allowbreak{}gather()},
        avec isolation des erreurs par coroutine. Si un scanner échoue, le pipeline
        continue avec les autres. Chaque résultat est mis en cache dans Redis avec un TTL
        d'une heure. L'\texttt{analyze\_\allowbreak{}review\_\allowbreak{}node} est le cœur du pipeline. Il envoie
        au modèle 14B le diff annoté avec numéros de lignes, les résultats nettoyés des
        scanners et le prompt combiné. Il extrait ensuite par regex le JSON tail contenant
        le score de risque, le verdict et les commentaires inline. Le \texttt{report\_node}
        finalise le traitement en publiant la revue sur GitHub, en fixant le statut du
        commit et en sauvegardant un fichier \texttt{summary.\allowbreak{}json} dans les artefacts.

        Le tableau \ref{tab:scan_matrix} présente la matrice de sélection des scanners selon
        la classification de la \textit{Pull Request}. Trivy FS et Gitleaks s'exécutent
        systématiquement pour toutes les catégories sauf docs. Les scanners additionnels
        sont activés selon la classification.

        \begin{table}[H]
        \centering
        \begin{tabular}{|m{3cm}|c|c|c|c|c|}
        \hline
        \textbf{Classification} & \textbf{Trivy FS} & \textbf{Gitleaks} & \textbf{Semgrep} & \textbf{Checkov} & \textbf{OSV} \\
        \hline
        \texttt{feature} & \checkmark & \checkmark & \checkmark & & \\
        \hline
        \texttt{dependency} & \checkmark & \checkmark & & & \checkmark \\
        \hline
        \texttt{infrastructure} & \checkmark & \checkmark & & \checkmark & \\
        \hline
        \texttt{config} & \checkmark & \checkmark & & & \\
        \hline
        \texttt{docs} & & & & & \\
        \hline
        \end{tabular}
        \caption{Matrice de sélection des scanners selon la classification de la PR}
        \label{tab:scan_matrix}
        \end{table}

        Lorsqu'un Dockerfile est détecté dans la \textit{Pull Request}, le nœud
        \texttt{scan\_\allowbreak{}full\_\allowbreak{}node} est activé à la place de \texttt{scan\_\allowbreak{}fs\_\allowbreak{}node}. Dans
        ce cas, l'image Docker est construite puis scannée par Trivy en complément du scan
        du filesystem.

        La catégorie \texttt{config} ne déclenche pour sa part ni Checkov ni OSV-Scanner.
        Checkov est réservé à la classification \texttt{infrastructure}, qui couvre les
        Dockerfile et les manifestes Terraform ou Kubernetes, alors qu'un simple fichier de
        configuration applicatif ne relève pas de l'Infrastructure-as-Code ; Trivy FS et
        Gitleaks y suffisent à couvrir respectivement les CVE de filesystem et les secrets
        éventuels.

        Deux fonctions de routage guident le flux du graphe. La fonction \texttt{route\_scans}
        redirige vers \texttt{skip\_\allowbreak{}scan\_\allowbreak{}node} pour les PRs de documentation, vers
        \texttt{scan\_\allowbreak{}full\_\allowbreak{}node} si un Dockerfile est détecté, et vers \texttt{scan\_\allowbreak{}fs\_\allowbreak{}node}
        dans tous les autres cas. La fonction \texttt{route\_risk} redirige vers
        \texttt{escalate\_\allowbreak{}node} lorsque le score de risque est CRITICAL ou HIGH et que
        l'escalade Slack est activée, et vers \texttt{report\_node} dans tous les autres cas.
        Tout nœud qui lève une exception positionne \texttt{state["error"]} et déclenche
        le routage vers \texttt{error\_node}.

        \subsubsection{Diagramme d'interaction — Réception du webhook}

        La figure \ref{fig:inter_webhook} présente le flux d'interaction entre GitHub,
        nginx, l'agent FastAPI et Redis lors de la réception d'un événement
        \textit{Pull Request}. L'agent valide d'abord la signature HMAC-SHA256, vérifie
        l'absence de doublon dans Redis, puis lance le pipeline LangGraph en tâche de fond
        et retourne immédiatement un code 202 à GitHub.

        \begin{figure}[H]
            \centering
            \frame{\includegraphics[width=0.88\columnwidth]{inter_webhook}}
            \caption{Diagramme d'interaction — Réception et validation du webhook GitHub}
            \label{fig:inter_webhook}
        \end{figure}

        \subsubsection{Diagramme d'interaction — Pipeline de revue}

        La figure \ref{fig:inter_pipeline} illustre la séquence d'interactions entre
        les nœuds LangGraph, Ollama, les outils SAST et l'API GitHub lors du traitement
        complet d'une \textit{Pull Request}. Le modèle 7B est d'abord interrogé pour la
        classification, puis les scanners s'exécutent en parallèle, et enfin le modèle
        14B produit la revue combinée qui est publiée sur GitHub.

        \begin{figure}[H]
            \centering
            \frame{\includegraphics[width=\textwidth]{inter_pipeline}}
            \caption{Diagramme d'interaction — Pipeline LangGraph de revue de sécurité}
            \label{fig:inter_pipeline}
        \end{figure}

    \subsection{Sous-système d'analyse du diff}

        \subsubsection{Architecture logique}

        Un modèle de langage peut halluciner des numéros de lignes qui n'existent pas dans
        le diff réel de la \textit{Pull Request}. Or, poster un commentaire \textit{inline}
        sur une ligne inexistante provoque une erreur 422 de l'API GitHub. Le module
        \texttt{diff\_\allowbreak{}parser.\allowbreak{}py} résout ce problème en construisant, pour chaque fichier
        modifié, l'ensemble exact des numéros de lignes valides avant toute publication.

        \begin{figure}[H]
            \centering
            \frame{\includegraphics[width=\columnwidth,height=0.78\textheight,keepaspectratio]{diff_parser_state_machine}}
            \caption{Machine à états du diff parser}
            \label{fig:diff_parser}
        \end{figure}

        Comme le montre la figure \ref{fig:diff_parser}, le parser parcourt le diff ligne
        par ligne et reconstruit, pour chaque fichier modifié, l'ensemble des numéros de
        lignes réellement présents dans la nouvelle version : les lignes ajoutées et les
        lignes de contexte alimentent cet ensemble, tandis que les lignes supprimées sont
        ignorées puisqu'elles n'existent pas dans le fichier final. Tout commentaire du
        modèle portant sur une ligne absente de cet ensemble est écarté avant publication,
        ce qui supprime les hallucinations de numéros de lignes.

        \subsubsection{Format du diff annoté}

        En plus du diff brut utilisé pour le récit de sécurité, un diff annoté de numéros
        de lignes visibles est transmis au modèle. Cette annotation lui permet de cibler
        précisément les lignes problématiques dans ses commentaires \textit{inline}. La
        figure \ref{fig:annotated_diff} illustre ce format.

        \begin{figure}[H]
            \centering
            \frame{\includegraphics[width=0.82\columnwidth]{annotated_diff_format}}
            \caption{Format du diff annoté avec numéros de lignes transmis au modèle 14B}
            \label{fig:annotated_diff}
        \end{figure}

    \subsection{Sous-système d'assistant de chat opérationnel}

        \subsubsection{Architecture logique}

        L'assistant de chat BTE Security AI Agent utilise une boucle ReAct
        (Reasoning + Acting) personnalisée. Cette conception est rendue nécessaire par le
        fait que \texttt{qwen2.\allowbreak{}5-coder} produit les appels d'outils en JSON texte brut
        et non via l'API native d'Ollama. La figure \ref{fig:react_loop_arch} illustre
        le flux complet de cette boucle.

        \begin{figure}[H]
            \centering
            \frame{\includegraphics[width=\columnwidth,height=0.80\textheight,keepaspectratio]{react_loop_architecture}}
            \caption{Boucle ReAct personnalisée}
            \label{fig:react_loop_arch}
        \end{figure}

        À chaque tour, les tokens produits par le modèle sont accumulés et analysés pour y
        détecter un appel d'outil au format JSON, qu'un extracteur tolérant sait reconnaître
        sous plusieurs formes (objet direct, bloc délimité, JSON embarqué dans du texte).
        Deux gardes encadrent la boucle : la première bloque tout appel identique déjà
        effectué dans la même réponse, ce qui supprime les boucles infinies ; la seconde
        force un appel d'outil dès qu'une question porte sur une donnée en temps réel (CPU,
        disque, conteneurs, alertes), pour empêcher le modèle d'y répondre de mémoire.
        L'outil est ensuite exécuté et son résultat injecté dans le contexte. La boucle se
        répète jusqu'à huit appels au maximum, après quoi la réponse finale est diffusée en
        événements SSE.

        Un modèle de langage a tendance à inventer des valeurs de métriques (pourcentage de
        CPU, espace disque libre, nombre de conteneurs) lorsqu'il répond à des questions sur
        l'état de l'infrastructure. Six couches de protection ont été conçues pour éliminer
        complètement ce comportement, comme présenté dans la figure
        \ref{fig:anti_hallucination_layers}.

        \begin{figure}[H]
            \centering
            \frame{\includegraphics[width=0.85\columnwidth]{anti_hallucination_layers}}
            \caption{Système anti-hallucination}
            \label{fig:anti_hallucination_layers}
        \end{figure}

        Les trois premières couches agissent sur la configuration du modèle :
        \texttt{temperature=0.\allowbreak{}0} rend la génération déterministe et supprime toute
        invention probabiliste de valeurs ; \texttt{num\_ctx} porté à 6 144 tokens laisse
        environ 4 100 tokens libres pour les observations des outils ; et
        \texttt{num\_predict} plafonné à 800 tokens garantit des réponses complètes. Les trois
        suivantes relèvent du contrôle applicatif : la garde sans-outil déjà évoquée, un bloc
        ANTI-HALLUCINATION de cinq règles dans le prompt système, et l'injection, après chaque
        appel d'outil, d'une observation imposant que toute valeur citée apparaisse verbatim
        dans un bloc \texttt{[OBSERVATION]}.

        L'assistant dispose de dix-neuf outils de monitoring couvrant l'ensemble de
        l'infrastructure. Le tableau \ref{tab:tools} présente ces outils organisés par
        catégorie.

        \begin{table}[H]
        \centering
        \begin{tabular}{|m{2.5cm}|m{5cm}|m{5.5cm}|}
        \hline
        \textbf{Catégorie} & \textbf{Outils} & \textbf{Données retournées} \\
        \hline
        VPS / Host & \texttt{vps\_status}, \texttt{disk\_usage}, \texttt{top\_\allowbreak{}processes}, \texttt{network\_\allowbreak{}stats}, \texttt{system\_\allowbreak{}net\_\allowbreak{}io} & CPU, RAM, disque, processus, interfaces réseau \\
        \hline
        Docker & \texttt{list\_\allowbreak{}containers}, \texttt{container\_\allowbreak{}logs}, \texttt{container\_\allowbreak{}stats}, \texttt{inspect\_\allowbreak{}container}, \texttt{list\_images}, \texttt{restart\_\allowbreak{}service} & État, logs, CPU/RAM live, santé \\
        \hline
        Ollama & \texttt{ollama\_\allowbreak{}status} & Modèles chargés, RAM consommée \\
        \hline
        Prometheus & \texttt{query\_\allowbreak{}prometheus}, \texttt{query\_\allowbreak{}prometheus\_\allowbreak{}range}, \texttt{prometheus\_\allowbreak{}alerts} & Métriques, tendances, alertes actives \\
        \hline
        Redis & \texttt{redis\_info} & Mémoire, hit rate, clients \\
        \hline
        Artifacts & \texttt{list\_\allowbreak{}scan\_\allowbreak{}artifacts}, \texttt{read\_\allowbreak{}scan\_\allowbreak{}artifact} & Fichiers SAST par PR \\
        \hline
        Base de données & \texttt{query\_\allowbreak{}database} & SELECT lecture seule \\
        \hline
        \end{tabular}
        \caption{Les 19 outils de monitoring du chat opérationnel}
        \label{tab:tools}
        \end{table}

        \subsubsection{Diagramme d'interaction — Boucle ReAct}

        La figure \ref{fig:inter_react} décrit le cycle d'interaction entre l'utilisateur,
        le serveur SSE, le modèle LLM et les outils de monitoring lors d'une question
        portant sur l'état de l'infrastructure. Le modèle génère un appel d'outil, l'agent
        l'exécute et injecte le résultat comme observation, puis le modèle produit la
        réponse finale en citant uniquement les valeurs présentes dans l'observation.

        \begin{figure}[H]
            \centering
            \frame{\includegraphics[width=0.88\columnwidth]{inter_react}}
            \caption{Diagramme d'interaction — Boucle ReAct du chat opérationnel}
            \label{fig:inter_react}
        \end{figure}

    \subsection{Sous-système de persistance}

        \subsubsection{Architecture logique}

        La base de données PostgreSQL comprend six tables applicatives, auxquelles s'ajoutent
        les quatre tables créées automatiquement par LangGraph pour la persistance des
        checkpoints. La figure \ref{fig:db_er} présente le diagramme entité-association
        de l'ensemble des tables applicatives.

        \begin{figure}[H]
            \centering
            \frame{\includegraphics[width=0.88\columnwidth]{database_er_diagram}}
            \caption{Diagramme E-A de la base de données PostgreSQL}
            \label{fig:db_er}
        \end{figure}

        \subsubsection{Description des tables}

        La table \texttt{pr\_reviews} est la table centrale de la base de connaissances.
        Elle conserve, pour chaque revue de \textit{Pull Request}, ses métadonnées (dépôt,
        numéro, auteur, classification), le score de risque et le verdict, le texte complet
        de la revue ainsi que le résumé des scanners et la liste des fichiers modifiés. La table
        \texttt{scan\_\allowbreak{}results} enregistre les sorties brutes de chaque scanner par dépôt et
        par déclencheur. La table \texttt{repo\_\allowbreak{}profiles} maintient un score de risque moyen
        glissant et un compteur de revues par dépôt, mis à jour après chaque exécution du
        nœud d'analyse. La table \texttt{security\_\allowbreak{}policies} contient les politiques de
        sécurité qui définissent les seuils de blocage, par exemple zéro CVE critique
        autorisé ou obligation de scan des secrets. La table \texttt{sbom\_cache} met en
        cache le Software Bill of Materials par dépôt pour éviter de rescanner les
        dépendances inchangées. Enfin, la table \texttt{incidents} permet de tracer les
        incidents de sécurité détectés via le chat ou les alertes.

        Les quatre tables LangGraph (\texttt{checkpoints}, \texttt{checkpoint\_\allowbreak{}blobs},
        \texttt{checkpoint\_\allowbreak{}writes} et \texttt{checkpoint\_\allowbreak{}migrations}) sont créées
        automatiquement par le checkpointer au démarrage de l'application. Elles contiennent
        l'état complet du graphe sérialisé et permettent la reprise automatique du pipeline
        après tout redémarrage du conteneur.

    \subsection{Sous-système d'observabilité}

        \subsubsection{Architecture logique}

        La stack d'observabilité centralise la surveillance de la plateforme. La figure
        \ref{fig:monitoring_stack} présente le flux de métriques depuis les sources de collecte
        jusqu'aux dashboards Grafana et aux alertes Slack.

        \begin{figure}[H]
            \centering
            \frame{\includegraphics[width=0.88\columnwidth]{monitoring_stack}}
            \caption{Architecture de la stack d'observabilité}
            \label{fig:monitoring_stack}
        \end{figure}

        L'agent expose vingt-huit métriques Prometheus personnalisées à l'endpoint
        \texttt{/metrics}, en plus des métriques HTTP auto-instrumentées. Treize couvrent
        l'agent, Ollama et le disque ; les quinze autres, ajoutées lors de l'extension
        d'observabilité décrite au chapitre suivant, se répartissent en trois familles :
        des métriques par conteneur (mémoire, CPU, trafic réseau, état), émises par un
        \textit{poller} interrogeant le socket Docker toutes les trente secondes ; des
        métriques de santé pour la sandbox LocalAI ; et une télémétrie A/B du routeur de
        chat, labellisée par couple \texttt{backend} et \texttt{model} afin de comparer
        Ollama et LocalAI sur les mêmes requêtes.

        La métrique \texttt{ollama\_\allowbreak{}reachable}, mise à jour toutes les trente secondes par
        un \textit{poller} asynchrone, distingue deux états qu'une simple jauge de modèles
        chargés confond : Ollama inactif entre deux revues, qui est l'état normal, et Ollama
        réellement injoignable. Cette distinction évite les fausses alertes critiques pendant
        les périodes de veille.

        Les quinze règles d'alerte sont organisées en cinq groupes : \textit{disk} (quatre
        règles sur l'occupation disque, côté node-exporter et côté agent), \textit{host}
        (trois règles : CPU, mémoire, I/O disque), \textit{agent} (trois règles :
        disponibilité, taux d'erreur, \textit{backlog} de revues), \textit{ollama} (deux
        règles : joignabilité et durée sans modèle chargé) et \textit{localai} (trois règles
        \texttt{info} sur la sandbox optionnelle). Une règle d'inhibition dans AlertManager
        supprime les alertes LocalAI redondantes lorsque la sandbox est arrêtée volontairement.

\section*{Conclusion}
\addcontentsline{toc}{section}{Conclusion}

Ce chapitre a couvert la conception complète du BTE Security AI Agent. L'architecture
globale repose sur onze conteneurs Docker organisés en quatre couches, exposés derrière
un unique point d'entrée nginx. La conception logicielle a ensuite été détaillée
sous-système par sous-système : l'inférence LLM à deux modèles, le pipeline LangGraph
à neuf nœuds piloté par des arêtes conditionnelles selon la classification de la PR,
le parser de diff qui élimine les hallucinations de numéros de lignes, l'assistant de
chat ReAct protégé par six couches anti-hallucination et ses dix-neuf outils de
monitoring, la base PostgreSQL à six tables applicatives, et la stack d'observabilité
étendue à vingt-huit métriques personnalisées. Chaque sous-système a suivi le même
plan : architecture logique, puis diagramme d'interaction. La concrétisation de ces
décisions, sprint par sprint, fait l'objet du chapitre suivant.
